\section*{Capítulo \ref{ch:b1:derivative} --- Derivación}

\begin{solution}{exo:b1:derivative:1}
$x^x = \eu^{x\ln x}$ en $\intoo{0}{+\infty}$: \omterm{def:b1:derivative:def}{derivada} $(\ln x +
1)\,x^x$.

$\ln(x + \sqrt{x^2+1})$ en $\R$ (el argumento es siempre $> 0$):
\omterm{def:b1:derivative:def}{derivada} $\frac{1}{\sqrt{x^2+1}}$ (calculada en
la~\cref{prop:b1:functions:invhyp} --- se trata de
$\operatorname{arsinh}$).

$\arctan\frac1x$ en $\R^*$: \omterm{def:b1:derivative:def}{derivada}
$\frac{-1/x^2}{1 + 1/x^2} = \frac{-1}{1 + x^2}$ (en consonancia con
la~\cref{prop:b1:functions:arcidentities}~(2): la función es
$\pm\frac\pi2 - \arctan x$ en cada semirrecta).

$\sqrt{1 + \eu^{2x}}$ en $\R$: \omterm{def:b1:derivative:def}{derivada}
$\frac{\eu^{2x}}{\sqrt{1 + \eu^{2x}}}$.
\end{solution}

\begin{solution}{exo:b1:derivative:2}
En $0$: $\bigl|\frac{f(h) - 0}{h}\bigr| = \abs{h \sin\frac1h} \leq
\abs h \to 0$, luego $f'(0) = 0$. Para $x \neq 0$, las reglas
habituales dan $f'(x) = 2x\sin\frac1x - \cos\frac1x$. A lo largo de
$x_n = \frac{1}{2\pi n}$: $f'(x_n) = 0 - 1 \to -1$; y a lo largo de
$y_n = \frac{1}{(2n+1)\pi}$: $f'(y_n) = 0 + 1 \to 1$. Dos sucesiones
que tienden a $0$ con límites distintos de $f'$: no hay límite
(\cref{thm:b1:continuity:seqchar}), luego $f'$ no es \omterm{def:b1:continuity:continuous}{continua} en $0$
y $f$ es \omterm{def:b1:derivative:def}{derivable} sin ser de clase $C^1$.
\end{solution}

\begin{solution}{exo:b1:derivative:3}
$\ln(1+x) < x$ para $x > 0$: desigualdad de tangencia por concavidad
en $0$ (estricta fuera del punto de contacto, pues $\ln$ es
estrictamente cóncava; o aplíquese el teorema del valor medio:
$\ln(1+x) = \frac{x}{1+c}$ para algún $c \in \intoo{0}{x}$, y
$\frac{x}{1+c} < x$). La misma identidad del valor medio da la cota
inferior: $\frac{x}{1+c} > \frac{x}{1+x}$.

Consecuencia: con $x/n$ en lugar de $x$,
\[
\frac{x/n}{1 + x/n} < \ln\Bigl(1 + \frac xn\Bigr) < \frac xn
\quad\implies\quad
\frac{x}{1 + x/n} < n \ln\Bigl(1 + \frac xn\Bigr) < x .
\]
El miembro izquierdo tiende a $x$: por emparedado,
$n\ln(1 + \frac xn) \to x$ y, por \omterm{def:b1:continuity:continuous}{continuidad} de $\exp$,
$\bigl(1 + \frac xn\bigr)^n = \eu^{n\ln(1 + x/n)} \to \eu^x$.
\end{solution}

\begin{solution}{exo:b1:derivative:4}
Sean $x_1 < x_2 < \dots < x_k$ raíces distintas de $P$. En cada
$\intcc{x_i}{x_{i+1}}$, Rolle (\cref{thm:b1:derivative:rolle})
produce $c_i \in \intoo{x_i}{x_{i+1}}$ con $P'(c_i) = 0$: son
$k - 1$ raíces de $P'$, distintas porque los \omterm{prop:b1:reals:intervals}{intervalos} \omterm{def:b1:topology:open}{abiertos} son
disjuntos --- y entrelazadas por construcción.

Si $P$ (de grado $n$) tiene todas sus raíces reales, escríbanse con
sus \omterm{def:b1:poly:derivative}{multiplicidades} $m_1 + \dots + m_k = n$. Cada raíz de
\omterm{def:b1:poly:derivative}{multiplicidad} $m_i \geq 2$ es raíz de $P'$ de \omterm{def:b1:poly:derivative}{multiplicidad}
$m_i - 1$ (\cref{prop:b1:poly:multiplicity}), lo que aporta
$\sum (m_i - 1) = n - k$; Rolle aporta $k - 1$ más, todas distintas
de las anteriores. Total $\geq n - 1 = \deg P'$: todas las raíces de
$P'$ son reales.
\end{solution}

\begin{solution}{exo:b1:derivative:5}
Fíjense $\varepsilon > 0$ y $A$ con
$\abs{f'(t) - \ell} \leq \varepsilon$ para $t \geq A$. Para $x > A$,
el teorema del valor medio en $\intcc{A}{x}$ da $c \in \intoo{A}{x}$
con
\[
f(x) = f(A) + f'(c)(x - A),
\qquad\text{luego}\qquad
\Bigl|\frac{f(x)}{x} - \ell\Bigr|
\leq \frac{\abs{f(A)} + \abs\ell A}{x} + \abs{f'(c) - \ell}
\cdot\frac{x - A}{x} \leq \frac{C_A}{x} + \varepsilon .
\]
Para $x$ grande, $\frac{C_A}{x} \leq \varepsilon$: por tanto,
$\frac{f(x)}{x} \to \ell$.

Sí: $f(x+1) - f(x) = f'(c_x)$ con $c_x \in \intoo{x}{x+1}$ (teorema
del valor medio en $\intcc{x}{x+1}$), y $c_x \to +\infty$, luego
$f(x+1) - f(x) \to \ell$.
\end{solution}

\begin{solution}{exo:b1:derivative:6}
Inducción sobre $n$. Para $n = 1$: Rolle. Si la afirmación vale para
$n - 1$: $f$ se anula en $n+1$ puntos, luego por Rolle aplicado en
los $n$ huecos, $f'$ se anula en $n$ puntos distintos; y la hipótesis
de inducción aplicada a $f'$ ($n-1$ veces \omterm{def:b1:derivative:def}{derivable}, con $n$ ceros)
hace que $(f')^{(n-1)} = f^{(n)}$ se anule en algún punto.

Aplicación: si $P$, de grado $\leq n$, se anula en $n+1$ puntos,
entonces $P^{(n)}$, que es una constante igual a $n!$ por el
coeficiente director, se anula: el coeficiente director es $0$, y se
concluye por inducción descendente (o directamente: todos los
coeficientes se anulan).
\end{solution}

\begin{solution}{exo:b1:derivative:7}
Por el teorema del valor medio en $\intcc{a}{x_0}$ y en
$\intcc{x_0}{b}$:
\[
f'(c_1) = \frac{f(x_0) - f(a)}{x_0 - a} = \frac{f(x_0)}{x_0 - a} > 0,
\qquad
f'(c_2) = \frac{f(b) - f(x_0)}{b - x_0} = \frac{-f(x_0)}{b - x_0} < 0,
\]
con $c_1 < x_0 < c_2$. Entonces el teorema del valor medio aplicado a
$f'$ en $\intcc{c_1}{c_2}$ da $c$ con
\[
f''(c) = \frac{f'(c_2) - f'(c_1)}{c_2 - c_1} < 0 . \qedhere
\]
\end{solution}

\begin{solution}{exo:b1:derivative:8}
$f'(x) = \frac{1 - \ln x}{x^2}$: $f$ crece en $\intoc{0}{\eu}$ y
decrece en $\intco{\eu}{+\infty}$, con máximo
$f(\eu) = \frac1\eu$; límites $-\infty$ en $0^+$ y $0$ en $+\infty$
(\omterm{prop:b1:functions:powerrules}{comparación de crecimientos}).

Para $\eu \leq a < b$: ser $f$ estrictamente decreciente ahí da
$\frac{\ln a}{a} > \frac{\ln b}{b}$, es decir, $b \ln a > a \ln b$,
es decir, $a^b > b^a$.

Con $a = \eu < b = \pi$: $\eu^\pi > \pi^\eu$.

Parejas pequeñas: $(2, 3)$: $2^3 = 8 < 9 = 3^2$ --- ¡al revés! La
razón: $2 < \eu$, y en $\intoo{0}{\eu}$ la función $f$ es
\emph{creciente}, de modo que la comparación se invierte cuando los
dos números están por debajo de $\eu$, y es imprevisible a caballo de
$\eu$ ($f(2) = f(4)$ explica el empate $2^4 = 4^2 = 16$).
\end{solution}

\begin{solution}{exo:b1:derivative:9}
\emph{Jensen para $\ln$, por inducción sobre $n$.} Afirmación: para
$x_i$ positivos y pesos $\lambda_i > 0$ con $\sum \lambda_i = 1$:
$\ln\bigl(\sum \lambda_i x_i\bigr) \geq \sum \lambda_i \ln x_i$. Para
$n = 2$ esto es la concavidad. Paso: con
$\Lambda = \lambda_1 + \dots + \lambda_{n-1} = 1 - \lambda_n$ y
$y = \sum_{i<n} \frac{\lambda_i}{\Lambda} x_i$,
\[
\ln\Bigl(\sum_{i \leq n} \lambda_i x_i\Bigr)
= \ln\bigl(\Lambda y + \lambda_n x_n\bigr)
\geq \Lambda \ln y + \lambda_n \ln x_n
\geq \Lambda \sum_{i<n} \frac{\lambda_i}{\Lambda}\ln x_i
+ \lambda_n \ln x_n,
\]
usando la concavidad ($n = 2$) y después la hipótesis de inducción.

Con $\lambda_i = \frac 1n$ y $x_i = a_i$:
$\ln\frac{\sum a_i}{n} \geq \frac 1n \sum \ln a_i =
\ln\sqrt[n]{a_1\cdots a_n}$; tómense exponenciales. Igualdad: $\ln$
es \emph{estrictamente} cóncava ($\ln'' < 0$), de modo que la
igualdad en cada paso obliga a que los puntos promediados coincidan
---es decir, a que todos los $a_i$ sean iguales---; y si todos son
iguales, la igualdad es clara.
\end{solution}

\begin{solution}{exo:b1:derivative:10}
Sea $g(x) = f(x) - vx$: \omterm{def:b1:derivative:def}{derivable}, con $g'(a) = f'(a) - v < 0$ y
$g'(b) = f'(b) - v > 0$. Por el teorema de los valores extremos, $g$
alcanza su mínimo en $\intcc{a}{b}$ en cierto $c$. No es en $a$:
como $g'(a) < 0$, los puntos justo a la derecha de $a$ tienen
$g < g(a)$. Tampoco es en $b$: como $g'(b) > 0$, los puntos justo a
la izquierda de $b$ tienen $g < g(b)$. Luego $c$ es \omterm{def:b1:topology:closure}{interior}, y
la~\cref{prop:b1:derivative:fermat} da $g'(c) = 0$, es decir,
$f'(c) = v$.
\end{solution}

\begin{solution}{exo:b1:derivative:11}
\emph{Unicidad:} dos puntos fijos $\ell \neq \ell'$ darían
$\abs{\ell - \ell'} = \abs{f(\ell) - f(\ell')} \leq k\abs{\ell -
\ell'} < \abs{\ell - \ell'}$, absurdo.

\emph{Existencia:} $g(x) = f(x) - x$ cumple, por la desigualdad del
valor medio, $f(x) \leq f(0) + k\abs x$; así, para
$x \geq \frac{\abs{f(0)}}{1 - k}$ se tiene
$g(x) \leq f(0) + kx - x \leq 0$, y simétricamente $g(-x) \geq 0$
para $x$ grande. El teorema del valor intermedio da un cero $\ell$ de
$g$: un punto fijo.

\emph{Convergencia:} otra vez la desigualdad del valor medio:
\[
\abs{u_{n+1} - \ell} = \abs{f(u_n) - f(\ell)} \leq k\abs{u_n - \ell},
\]
luego, por inducción, $\abs{u_n - \ell} \leq k^n \abs{u_0 - \ell} \to 0$.
\end{solution}

\begin{solution}{exo:b1:derivative:12}
\begin{enumerate}
  \item Si $g(b) = g(a)$, Rolle daría un cero \omterm{def:b1:topology:closure}{interior} de $g'$:
        excluido. Póngase $\lambda = \frac{f(b) - f(a)}{g(b) -
        g(a)}$ y $h = f - \lambda g$: $h$ es \omterm{def:b1:continuity:continuous}{continua} en
        $\intcc{a}{b}$, \omterm{def:b1:derivative:def}{derivable} dentro, y $h(b) - h(a)
        = f(b) - f(a) - \lambda(g(b) - g(a)) = 0$. Rolle
        proporciona $c$ con $h'(c) = 0$, es decir, $f'(c) =
        \lambda\,g'(c)$; divídase por $g'(c) \neq 0$.
  \item Para $x > a$ próximo a $a$, el apartado (1) en
        $\intcc{a}{x}$ (donde $g' \neq 0$) da $g(x) \neq 0$ y
        $c_x \in \intoo{a}{x}$ con
        \[
        \frac{f(x)}{g(x)} = \frac{f(x) - f(a)}{g(x) - g(a)}
        = \frac{f'(c_x)}{g'(c_x)} .
        \]
        Cuando $x \to a^+$, $c_x \to a^+$ (emparedado), luego el
        miembro derecho tiende a $\ell$: $\frac{f}{g} \to \ell$.
  \item $\frac{f(x)}{g(x)} = x\sin\frac1x \to 0$, mientras que
        $\frac{f'(x)}{g'(x)} = 2x\sin\frac1x - \cos\frac1x$ no
        tiene límite en $0$ (\cref{exo:b1:derivative:2}): la regla
        de l'Hôpital transfiere información solo de
        $\frac{f'}{g'}$ a $\frac fg$, nunca al revés.
\end{enumerate}
\end{solution}

\begin{solution}{pb:b1:derivative:1}
\textbf{1.} $\bigl|\frac ab - \frac pq\bigr| = \frac{\abs{aq -
bp}}{bq}$, y $aq - bp$ es un entero no nulo cuando las fracciones
son distintas: la distancia es $\geq \frac{1}{bq}$. Así pues, ningún
racional distinto del propio $x$ entra en el \omterm{prop:b1:reals:intervals}{intervalo} perforado de
radio $\frac{1}{bq}$ en torno a $x = \frac ab$.

\textbf{2.} Si $\bigl|\sqrt2 - \frac pq\bigr| \geq 1 \geq
\frac{1}{4q^2}$, listo. En caso contrario,
$\frac pq \in \intoo{\sqrt2 - 1}{\sqrt2 + 1}$, luego
$0 < \sqrt2 + \frac pq < 2\sqrt2 + 1 < 4$. Como $\sqrt 2 \notin \Q$,
$p^2 - 2q^2$ es un entero no nulo, y
\[
\Bigl|\sqrt2 - \frac pq\Bigr|
= \frac{\abs{2q^2 - p^2}}{q^2\,\bigl(\sqrt2 + \frac pq\bigr)}
\geq \frac{1}{4q^2} .
\]

\textbf{3.} $(p + 2q)^2 - 2(p + q)^2 = -(p^2 - 2q^2)$: el valor
$\pm1$ se propaga. A partir de $(1,1)$:
\[
(3,2),\ (7,5),\ (17,12),\ (41,29),\ (99,70),
\]
con $p^2 - 2q^2$ alternando $-1, +1, \dots$
Para estas, $\frac pq \geq 1$, luego $\sqrt2 + \frac pq > 2$ y
\[
\Bigl|\sqrt2 - \frac pq\Bigr| =
\frac{1}{q^2(\sqrt2 + p/q)} < \frac{1}{2q^2} ,
\]
con $q \to \infty$: infinitas aproximaciones de orden $2$. Junto con
la pregunta 2, el exponente $2$ es exacto para $\sqrt 2$.

\textbf{4.} Los $N + 1$ números $kx - \lfloor kx\rfloor$
($0 \leq k \leq N$) están en las $N$ cajas
$\intco{\frac jN}{\frac{j+1}{N}}$: dos comparten caja
(\cref{cor:b1:counting:pigeonhole}), digamos para $i < j$. Con
$q = j - i \leq N$ y $p = \lfloor jx\rfloor - \lfloor ix\rfloor$:
$\abs{qx - p} < \frac1N$, luego
$\bigl|x - \frac pq\bigr| < \frac{1}{Nq} \leq \frac{1}{q^2}$.
Haciendo $N \to \infty$: como $x$ es irracional, cada fracción fija
está a distancia positiva de $x$, mientras que
$\frac{1}{Nq} \leq \frac 1N \to 0$ obliga a que aparezcan fracciones
nuevas: infinitas $\frac pq$ distintas con
$\bigl|x - \frac pq\bigr| < \frac{1}{q^2}$.

\textbf{5.} Pártase de cualquier $P_0$ entero no nulo con
$P_0(x) = 0$. Si $P_0$ tiene una raíz racional $\frac ab$, el teorema
del factor (\cref{thm:b1:poly:factor}) escribe
$P_0 = \bigl(X - \frac ab\bigr)Q$ con $Q \in \Q[X]$; como
$x \neq \frac ab$ ($x$ es irracional), $Q(x) = 0$, y quitando
denominadores se obtiene un \omterm{def:b1:poly:def}{polinomio} \emph{entero} no nulo de grado
menor que se anula en $x$. El grado baja en cada paso, así que el
proceso termina: se llega a $P \in \Z[X]$, $P(x) = 0$, sin raíz
racional, de cierto grado $d$. Si $d \leq 1$, $P = uX + v$ haría
racional a $x = -\frac vu$: luego $d \geq 2$.

\textbf{6.} $q^d\,P\bigl(\frac pq\bigr) = a_d p^d + a_{d-1}
p^{d-1} q + \dots + a_0 q^d$ es un entero, y es no nulo porque $P$ no
tiene raíz racional: $\bigl|P\bigl(\frac pq\bigr)\bigr| \geq q^{-d}$.

\textbf{7.} Nótese que $M > 0$: $P'$ es un \omterm{def:b1:poly:def}{polinomio} no nulo
($d \geq 2$), luego no puede anularse idénticamente en
$\intcc{x-1}{x+1}$. Si $\bigl|x - \frac pq\bigr| > 1$, entonces
supera $\frac{C}{q^d}$ trivialmente. En caso contrario,
$\frac pq \in \intcc{x-1}{x+1}$ y el teorema del valor medio
(\cref{thm:b1:derivative:mvt}) da $c$ entre $x$ y $\frac pq$ con
\[
\Bigl|P\Bigl(\frac pq\Bigr)\Bigr|
= \Bigl|P\Bigl(\frac pq\Bigr) - P(x)\Bigr|
= \abs{P'(c)}\,\Bigl|x - \frac pq\Bigr|
\leq M\,\Bigl|x - \frac pq\Bigr| ,
\]
luego, con la pregunta 6: $\bigl|x - \frac pq\bigr| \geq
\frac{1}{Mq^d} \geq \frac{C}{q^d}$.

\textbf{8.} Supóngase que $x = \frac ab$ es de Liouville. Tómense $n$
con $2^{n-1} > b$ y el $\frac pq$ correspondiente, $q \geq 2$:
\[
0 < \Bigl|x - \frac pq\Bigr| < \frac{1}{q^n}
= \frac{1}{q^{n-1}\,q} \leq \frac{1}{2^{n-1} q} <
\frac{1}{bq} ,
\]
en contradicción con la pregunta 1. Luego los números de Liouville son irracionales.

\textbf{9.} Si un $x$ de Liouville fuese algebraico: es irracional
(pregunta 8), luego las preguntas 5--7 proporcionan $d \geq 2$ y
$C > 0$ con $\bigl|x - \frac pq\bigr| \geq \frac{C}{q^d}$ siempre.
Para cada $n$, el aproximante de Liouville da
$\frac{C}{q^d} < q^{-n}$, es decir, $C < q^{d-n} \leq 2^{d-n}$ (pues
$q \geq 2$). Y para $n$ grande, $2^{d-n} < C$: contradicción. Los
números de Liouville son trascendentes.

\textbf{10.} Unos en las posiciones $1, 2, 6, 24$; todas las demás
cifras entre las $25$ primeras se anulan:
\[
L = 0.1100010000\,0000000000\,00010\dots
\]

\textbf{11.} Para $m > k$, las posiciones $n!$ con $n > k$ son
enteros distintos $\geq (k+1)!$, de modo que la suma geométrica
finita da
\[
t_m - t_k = \sum_{n=k+1}^{m} 10^{-n!}
\leq \sum_{j = (k+1)!}^{m!} 10^{-j}
< 10^{-(k+1)!}\,\frac{1}{1 - \frac1{10}}
= \frac{10}{9}\,10^{-(k+1)!} ;
\]
tomando el \omterm{def:b1:reals:bounds}{supremo} sobre $m$:
$L - t_k \leq \frac{10}{9}10^{-(k+1)!} < 2\cdot10^{-(k+1)!}$. Cota
inferior: $L \geq t_{k+1} = t_k + 10^{-(k+1)!}$.

\textbf{12.} $p_k = 10^{k!}\,t_k \in \N$, $q_k = 10^{k!}$, y
$(k+1)! = (k+1)\,k!$ da $10^{-(k+1)!} = q_k^{-(k+1)}$: la
pregunta 11 dice
\[
0 < L - \frac{p_k}{q_k} < \frac{2}{q_k^{\,k+1}} .
\]
Dado $n$: para $k \geq n$, $2\,q_k^{-(k+1)} \leq q_k^{-n}$ (en
efecto, $q_k^{\,k+1-n} \geq q_k \geq 10 > 2$), y $q_k \geq 2$: se
cumple la definición de la pregunta 8. $L$ es un número de Liouville.

\textbf{13.} Por la pregunta 9, $L$ es trascendente --- el primer
número de la historia del que se demostró la trascendencia
(Liouville, 1844). Comprobación con las cifras: la cadena tiene
infinitos unos con huecos consecutivos
$(k+1)! - k! = k\cdot k! \to \infty$, luego no es finalmente
periódica, y $L \notin \Q$ por el criterio de periodicidad
del~\cref{pb:b1:reals:1} (pregunta 18) --- coherente.

\textbf{14.} Con cifras $d_n \in \intint{1}{9}$ en las posiciones
factoriales: la cota de la cola de la pregunta 11 se multiplica a lo
sumo por $9$: $0 < L' - t'_k \leq
9\cdot\frac{10}{9}\,10^{-(k+1)!} = 10\,q_k^{-(k+1)}$ (la positividad,
porque la cifra en la posición $(k+1)!$ es no nula). Para $k \geq n$:
$10\,q_k^{-(k+1)} \leq q_k^{-n}$, pues $q_k^{\,k+1-n} \geq 10$: de
nuevo un número de Liouville y, por tanto, trascendente. Estos
valores son distintos dos a dos para elecciones de cifras distintas
(las cadenas son propias ---abundan los ceros--- y las cadenas
propias determinan su valor, \cref{pb:b1:reals:1}, pregunta 10). Dada
una lista cualquiera $k \mapsto x_k$ de ellos, elíjase la $k$-ésima
cifra factorial en $\intint{1}{9}$ distinta de la de $x_k$: un número
de la misma forma que falta en la lista. No numerablemente muchos
trascendentes explícitos.

\textbf{15.} Primero un lema: \emph{si
$\bigl|x - \frac pq\bigr| \geq \frac{C}{q^s}$ para todo
$\frac pq \neq x$, entonces $x$ no es aproximable a ningún orden
$\mu > s$.} En efecto, infinitas $\frac pq \neq x$ con
$\bigl|x - \frac pq\bigr| < \frac{c}{q^\mu}$ forzarían
$\frac{C}{q^s} < \frac{c}{q^\mu}$, es decir,
$q^{\mu - s} < \frac cC$: los $q$ están acotados, y solo hay una
cantidad acotada de fracciones a distancia $1$ de $x$ --- finitos
candidatos, no infinitos. Ahora la jerarquía: los racionales son
aproximables al orden $1$ ($\frac pq$ con
$p = \lfloor qx\rfloor + 1$ da error $\leq \frac1q < \frac2q$) y a
ningún orden $\mu > 1$ (la pregunta 1 da la hipótesis del lema con
$s = 1$, $C = \frac1b$); $\sqrt2$: al orden $2$ (pregunta 3) y no
más (pregunta 2 y el lema); todo irracional: al menos $2$
(pregunta 4); los algebraicos de grado $d$: a lo sumo $d$
(pregunta 7 y el lema); y los números de Liouville: a todo orden (la
fórmula de la pregunta 12, con $c = 2$).

\textbf{16.} Con $r = \frac ab$:
$\frac{p_k}{q_k} + \frac ab = \frac{b p_k + a q_k}{b q_k} =:
\frac{P_k}{Q_k}$, $Q_k = b q_k \geq 2$, y
\[
\Bigl|(L + r) - \frac{P_k}{Q_k}\Bigr| = L - \frac{p_k}{q_k}
< 2\,q_k^{-(k+1)} = 2\,b^{\,k+1} Q_k^{-(k+1)} .
\]
Dado $n$: para $k$ grande,
$Q_k^{\,k+1-n} \geq Q_k = b\,10^{k!} \geq 2\,b^{\,k+1}$ (el factorial
aplasta a la potencia), luego el error es $< Q_k^{-n}$: $L + r$ es de
Liouville. Y como $\Q$ es \omterm{def:b1:topology:dense}{denso} y cada $L + r$ es trascendente, los
\omterm{pb:b1:logic:1}{números trascendentes} son \omterm{def:b1:topology:dense}{densos} en $\R$.

\textbf{17.} Para $h \geq 1$ hay finitos $P \in \Z[X]$ con
$\deg P + \sum_i \abs{a_i} \leq h$ (grado $\leq h$ y cada
coeficiente en $\intint{-h}{h}$: a lo sumo $(2h+1)^{h+1}$). Todo
\omterm{def:b1:poly:def}{polinomio} entero no nulo tiene una altura así, y tiene a lo sumo
$\deg P$ raíces reales: los \omterm{pb:b1:logic:1}{números algebraicos} forman una unión
numerable (sobre $h$) de \omterm{def:b1:counting:card}{conjuntos finitos} y, por tanto, se pueden
enumerar en una sola sucesión. Si los trascendentes también se
pudieran enumerar, intercalar las dos listas enumeraría $\R$, en
contradicción con el~\cref{pb:b1:reals:1} (pregunta 22). Luego los
\omterm{pb:b1:logic:1}{números trascendentes} forman un \omterm{def:b1:logic:sets}{conjunto} no numerable.
Comparación: Cantor demuestra que \emph{casi todos} los reales son
trascendentes sin exhibir ninguno; Liouville exhibe uno, con
constantes efectivas (parte V) --- existencia por abundancia frente a
existencia por construcción.

\textbf{18.} Por la pregunta 2, la hipótesis del lema se cumple con
$s = 2$, $C = \frac14$. Si $\sqrt2$ fuera de Liouville, entonces para
$n = 3$: $\frac{1}{4q^2} < q^{-3}$ obliga a $q < 4$, luego
$q \in \{2, 3\}$; y solo hay finitas $\frac pq$ con esos $q$ a
distancia $1$ de $\sqrt2$, cada una a cierta distancia positiva
$\geq \varepsilon_0$ ($\sqrt2$ es irracional); eligiendo $n$ con
$2^{-n} < \varepsilon_0$ no queda ninguna $\frac pq$ admisible:
contradicción. El mismo argumento con $\frac{C}{q^d}$ muestra que
ningún \omterm{pb:b1:logic:1}{número algebraico} es de Liouville --- la pregunta 9 con ropa
efectiva.

\textbf{19.} $99^2 - 2\cdot70^2 = 9801 - 9800 = 1$. Por tanto,
\[
\sqrt2 - \frac{99}{70} =
\frac{2 - (99/70)^2}{\sqrt2 + 99/70}
= \frac{-1}{4900\,\bigl(\sqrt2 + \tfrac{99}{70}\bigr)} ,
\qquad
\Bigl|\sqrt2 - \frac{99}{70}\Bigr|
= \frac{1}{4900 \times 2.8284\dots} \approx 7.2\cdot10^{-5} :
\]
$\frac{99}{70} = 1.414285\dots$ frente a
$\sqrt2 = 1.414213\dots$ --- cinco cifras correctas.

\textbf{20.} Test de raíces racionales para $P = X^3 - 2$:
candidatas $\pm1, \pm2, \pm\frac12$, y ninguna es raíz. Luego $d = 3$
y la parte II se aplica a $x = 2^{1/3} = 1.2599\dots$ En
$\intcc{x - 1}{x + 1} \subseteq \intcc{0.25}{2.26}$:
$\abs{P'(t)} = 3t^2 \leq 3\,(1 + 2^{1/3})^2 < 3\times(2.26)^2 =
15.32 < 16$, luego $M < 16$ y $C \geq \frac{1}{16}$:
\[
\Bigl|2^{1/3} - \frac pq\Bigr| \geq \frac{1}{16\,q^3}
\qquad\text{para todos los racionales.}
\]

\textbf{21.} Si $\bigl|2^{1/3} - \frac pq\bigr| < 10^{-6}$, entonces
$\frac{1}{16 q^3} < 10^{-6}$, es decir,
$q^3 > \frac{10^6}{16} = 62\,500$; y como
$39^3 = 59\,319 < 62\,500 \leq 64\,000 = 40^3$: $q \geq 40$.

\textbf{22.} Ejecútese la parte III en base $2$:
$B = \sup_k \sum_{n\leq k} 2^{-n!}$, $q_k = 2^{k!}$, y la cola
geométrica (de razón $\frac12$) da
$0 < B - \frac{p_k}{q_k} \leq 2\cdot2^{-(k+1)!} = 2\,q_k^{-(k+1)}
\leq q_k^{-n}$ para $k \geq n$. Luego $B$ es de Liouville y, por
tanto, trascendente: nada en el argumento es decimal.

\textbf{23.} Con unos en las posiciones $3^k$: $q_k = 10^{3^k}$ y la
cota de la cola da
$0 < x^\dagger - \frac{p_k}{q_k} < 2\cdot10^{-3^{k+1}} =
2\,q_k^{-3}$ (pues $3^{k+1} = 3\cdot3^k$): infinitas aproximaciones
de orden $3$. Por el lema de la pregunta 15: el orden $3 > 1$ excluye
la racionalidad, y el orden $3 > 2$ excluye ser un irracional
cuadrático (cuya desigualdad de Liouville tiene $s = d = 2$). Pero un
\omterm{pb:b1:logic:1}{número algebraico} de grado $\geq 3$ solo es repelido al orden
$d \geq 3$: el método de Liouville no puede separar $x^\dagger$ de
los cúbicos. La brecha la cierra el teorema de Roth ---todo
irracional algebraico tiene orden de aproximación exactamente $2$---,
un resultado del siglo XX muy por encima de este volumen; admitido
este, $x^\dagger$ también es trascendente.

\textbf{24.} Hay a lo sumo $(2H+1)^{d+1}$ tuplas
$(a_0, \dots, a_d)$ con entradas en $\intint{-H}{H}$, y cada
\omterm{def:b1:poly:def}{polinomio} no nulo de entre ellos tiene a lo sumo $d$ raíces reales:
surgen a lo sumo $d\,(2H+1)^{d+1}$ \omterm{pb:b1:logic:1}{números algebraicos} --- la
finitud que permitió a la pregunta 17 enumerarlos todos.

\textbf{25.} (i) El único ingrediente analítico es el teorema del
valor medio, en la pregunta 7, que convierte la anulación
$P(x) = 0$ en la repulsión lipschitziana
$\abs{P(p/q)} \leq M\abs{x - p/q}$. (ii) La tensión: la
integralidad empuja $\abs{P(p/q)}$ hacia arriba, hasta $q^{-d}$, y la
suavidad tira de él hacia abajo, hasta $M\abs{x - p/q}$ --- un
racional demasiado próximo a $x$ quedaría aplastado entre ambas.
(iii) Liouville construye un trascendente con constantes efectivas;
Cantor demuestra que casi todos los reales son trascendentes sin
nombrar ninguno: construcción frente a \omterm{def:b1:counting:card}{cardinal}. (iv) El problema
del fin de semana del~\cref{ch:b1:integration} demuestra la
irracionalidad de $\pi$ con el mismo emparedado ---una integral que
sería un entero positivo y, sin embargo, queda atrapada en
$\intoo{0}{1}$--- con la integración en el papel analítico que aquí
hace la derivación.
\end{solution}
